% Generated draft: requires mathematical and source review.
\documentclass[11pt]{article}
\usepackage[margin=25mm]{geometry}
\usepackage{fontspec}
\setmainfont{Latin Modern Roman}
\usepackage{amsmath,amssymb,mathtools,unicode-math}
\setmathfont{Latin Modern Math}
\usepackage{enumitem,needspace,xurl,hyperref}
\setlength{\parindent}{0pt}
\setlength{\parskip}{5pt}
\setlength{\emergencystretch}{3em}
\newcommand{\N}{\mathbb{N}}
\newcommand{\Z}{\mathbb{Z}}
\newcommand{\Q}{\mathbb{Q}}
\newcommand{\R}{\mathbb{R}}
\newcommand{\C}{\mathbb{C}}
\newcommand{\F}{\mathbb{F}}
\newcommand{\A}{\mathbb{A}}
\newcommand{\PP}{\mathbb{P}}
\newcommand{\E}{\mathbb{E}}
\newcommand{\eps}{\varepsilon}
\newcommand{\dd}{\,\mathrm{d}}
\begin{document}
\section*{UnsolvedMath ID: 20001973}
\textbf{Source problem code:} catalogue entry
\subsection{Definitions and mathematical statement}
N. Proudfoot's Question 9.1 asks whether singular hyperk\"ahler quotients have a natural ring structure on intersection cohomology. It is motivated by a seemingly natural structure in singular hypertoric examples, described as a quotient of \(H_T^*(T^*\mathbb C^n)\). Comment 9.2 notes that a small resolution gives a ring via ordinary cohomology, while the author was then unaware of such resolutions in the hypertoric examples. This is a historical comment, not a theorem that such resolutions never exist.

For a singular quotient \(X\), use rational intersection cohomology with classical lower-middle perversity
\[
 \overline m(c)=\left\lfloor\frac{c-2}2\right\rfloor
\]
at occurring singular strata of real codimension \(c\geq2\). The desired internal product would have the form
\[
 IH_{\overline m}^i(X;\mathbb Q)\otimes
 IH_{\overline m}^j(X;\mathbb Q)
 \longrightarrow IH_{\overline m}^{i+j}(X;\mathbb Q),
\]
with a specified meaning of naturality and a unit. The source does not supply that meaning, coefficient conventions or a category of quotients.

The aid reports a later positive hypertoric result of Braden--Proudfoot: a natural algebra on equivariant hypertoric intersection cohomology, and for central unimodular arrangements a unique commutative unital ring object on the equivariant intersection complex. The precise categorical and normalization hypotheses are not reproduced. It does not assert that this construction extends formally to arbitrary singular hyperk\"ahler quotients. Likewise, transporting a ring from one small resolution does not by itself make it independent of every possible choice of resolution.

The aid's diagnostic concerns the standard Goresky--MacPherson general-position pairing. Its perversity bound can return to the same lower-middle target by the usual closure test only if
\[
 2\overline m(c)\leq\overline m(c)
 \quad\text{at every occurring singular codimension}.
\]
Since \(\overline m(c)\geq0\), this formal test is equivalent to \(\overline m(c)=0\) everywhere, or to absence of singular strata of real codimension at least four. It is a test for the unaugmented pairing construction, not a prohibition on other products with extra geometric or sheaf data.

At a real codimension-four stratum, \(\overline m(4)=1\). If two allowable chains have dimensions \(i,j\) in an ambient real dimension \(N\), the general-position bound for their intersection with that stratum is
\[
 (i-3)+(j-3)-(N-4)=i+j-N-2.
\]
A product chain of dimension \(k=i+j-N\) is only lower-middle allowable there up to dimension \(k-3\). Thus the generic dimension estimate exceeds the middle bound by one; top perversity \(\overline t(4)=2\) permits \(k-2\). This compares allowable bounds and does not claim every pair of actual cycles saturates the excess. The diagnostic explains why ordinary middle-perversity closure is not automatic, while leaving open the broader requested natural-ring construction.
\subsection{Short English statement}
The source asks for a natural intersection-cohomology ring on general singular hyperk\"ahler quotients. The aid reports a positive hypertoric construction and a limitation of the ordinary intersection pairing: its lower-middle perversity bound is not closed at codimension four. This diagnoses the need for additional structure rather than proving that every possible ring is impossible.
\subsection{Sources}
\begin{itemize}
\item Local source record: catalogue entry, catalogue entry; canonical \texttt{aim-geometry-notes.json}, \texttt{notes[310]}, source item 9.1. Title: ``A perversity-closure diagnostic for intersection-cohomology rings.''
\item AIM workshop list \emph{Moment maps and surjectivity in various geometries}, Question 9.1 and Comment 9.2 attributed to N. Proudfoot; supplied reference \url{https://aimath.org/WWN/momentmaps/momentmaps.pdf}. This draft uses the supplied text without accessing a PDF. Original remarks and literature fields are empty.
\item Supplied research attributions: Braden--Proudfoot for hypertoric equivariant intersection-cohomology algebras and the central-unimodular ring-object result; Goresky--MacPherson for the standard perversity pairing. No full bibliographic entries or exact categorical hypotheses are supplied. The explicit codimension-four diagnostic belongs to attempt 1.
\item Registry dated 14 August 2026: ``partially\_solved,'' ``PARTIAL-PROGRESS,'' statement status ``exact''; attempt partial, rigor ``proved,'' novelty confidence low. The machine-generated, non-peer-reviewed aid requires independent verification before citation.
\end{itemize}
\end{document}
